% Folder 34, batch 06 — archivists' pages 101–120.
% Pass: Opus 5.5 (claude-opus-5-5), 2026-10-03 — first pass, unchecked against the pages by a human.
%
\documentclass[11pt,a4paper]{article}
\input{../preamble/grothendieck}

\folder{34}
\batch{6}
\pages{101}{120}
\foldertitle{SGA 7 : notes manuscrites (s.d.).}
\dating{[vers 1967-1973]}
\watermark{Édition de démonstration}

\begin{document}

\section{4. Calculs \add{cohomologiques} locaux liés au théorème de pureté}

\page{101}
\note{Pagination de l'auteur : « 17 » ; un « III » souligné en tête du feuillet, à gauche. Le paragraphe 4 commence ici.}

\marginal{\uncertain{devrait commencer au n° 3}}

\textbf{Définition 4.1.} a) Soit $X$ un schéma régulier. On dit que $X$ \emph{satisfait au théorème de pureté cohomologique} si, pour tout \struck{\ill{}} $X'$ étale sur $X$, \ill{} \uncertain{tout} sous-schéma régulier $Y$ de $X'$, \struck{\ill{}} et tout sous-schéma régulier $Z$ de $Y$ purement de codimension $1$, \struck{$\underline{H}^i_Z($} et tout entier $n$ premier aux caractéristiques résiduelles de $Z$, on a (où $\Lambda = \mathbf{Z}/n\mathbf{Z}$)
\[
(4.1.1) \qquad \underline{H}^i_Z(\Lambda_Y) = 0 \quad \text{pour } i \neq 2
\]
ou, ce qui revient au même, si $f : U = Y - Z \to Y$ est l'inclusion canonique,
\[
(4.1.2) \qquad R^i f_*(\Lambda_U) = 0 \quad \text{pour } i \neq 0, 1 .
\]
b) \add{Si $r$ est un entier,} \struck{On dit que} $X$ satisfait au théorème de pureté cohomologique en dimension $\leq r$ si (4.1.2) est vrai pour $i \leq r$.

\textbf{Remarques 4.2.} a) La condition énoncée \add{dans 4.1 a)} \struck{\ill{}} est locale pour la top.\ (ét). Elle \struck{s'exp} s'exprime \add{aussi} en disant que pour \uncertain{tout} localisé strict $X'$ de $X$, \ill{} sous-schéma régulier $Y$ de $X'$ \add{et tt $n$ premier à la car.\ rés.\ de $X'$}, et tout diviseur régulier $Z$ de $Y$, on a (4.1.1) ou, ce qui revient au même,
\[
H^i(Z - Y, \mathbf{Z}/n\mathbf{Z}) = 0 \quad \text{si } i \neq 0, 1
\]
(resp.\ et $i \leq r$). \struck{\ill{}}\note{« $Z - Y$ » ainsi sur la page, là où l'on attend $Y - Z$.}

b) Il est évident \add{(simple définition)} que \struck{\ill{}} $X$ satisfait au th.\ de pureté \add{coh.} en dimension $\leq 1$.

\page{102}
\note{Pagination de l'auteur : « 18 ».}

c) Rappelons qu'on conjecture que tout schéma régulier satisfait au th.\ de pureté cohomologique (SGA 4 XVI 3.10 b). Cette propriété est démontrée par l'\uncertain{auteur} dans les cas suivants :

1°) $X$ est lisse sur un corps quelconque (plus général\supplied{ement}, la relation (4.1.1) ou (4.1.2) est vraie si $Y$, $Z$ sont lisses sur un corps quelconque), cf.\ (SGA 4 XVI 3.7, 3.9).

2°) $X$ est excellent de caractéristique nulle (SGA 4 XIX 3.2). \struck{Plus général}

\struck{Mais pour} Plus généralement, \add{en se limitant aux loc.\ strict.\ \uncertain{près},} si $X$ est excellent de car.\ donnée, et si on dispose de la résolution des singularités des schémas excellents en cette car., alors $X$ satisfait au th.\ de pureté \add{\uncertain{aucun}} \struck{\ill{}} \ill{}, (pour \ill{} schémas réguliers d'inégales caractéristiques et de dimension $\geq 2$, le théorème de pureté n'est \uncertain{prouvé} à l'heure actuelle \ill{} \ill{}.

d) Les conditions de 4.1 \uncertain{est stable} par localisation étale et par passage à un \ill{} \uncertain{sous-schéma régulier}.

\textbf{Proposition 4.2.} Supposons que $X$ satisfait au th.\ de pureté cohomologique, et soit $Y$ un sous-schéma fermé régulier de $X$ purement de codimension $d$. Alors on a un isomorphisme canonique, si $\Lambda = \mathbf{Z}/n\mathbf{Z}$ ($n$ premier aux car.\ rés.\ de $X$) :
\[
(4.2.1) \qquad \underline{H}^i_Y(\Lambda_X) = 0 \quad \text{si } i \neq 2d ,
\]
\note{après le signe $=$, une première rédaction biffée, partiellement lisible : $\mu_n^{\otimes(-d)}$ … si $i = 2d$.}
\[
(4.2.2) \qquad \underline{H}^{2d}_Y(\Lambda_X) \simeq \bigl(\mu_n^{\otimes(-d)}\bigr)_Y
\]
(ou encore
\[
(4.2.2 \text{ bis}) \qquad \underline{H}^{2d}_Y(\mu_n^{\otimes d}) \simeq \Lambda_Y\ ).
\]

\page{103}
\note{Pagination de l'auteur : « 19 ».}

\emph{Démonstration.} Prouvons d'abord (4.2.1) et le fait que $\underline{H}^{2d}_Y(\Lambda_X)$ est loc.\ isom.\ à $\Lambda_Y$ : \struck{récurrence sur $d$}\struck{(4.2.1), \ill{} question étant locale}, le cas $d = 0$ étant trivial, le cas $d = 1$ résulte de la théorie de Kummer et du lemme de Abhyankar absolu, \struck{\ill{}} qui fournit $\underline{H}^2_Y(\mu_{n\,X}) \simeq \Lambda_Y$ (cf.\ SGA 4 XIX 1.3 et SGA 1 XII …). Pour $d > 1$, on procède par récurrence sur $d$. La question étant locale, \struck{\ill{}} on peut supposer qu'il existe un diviseur \emph{régulier} $Z$ dans $X$ tel que $Y \subset Z$. \struck{On utilise} Considérons la suite spectrale de transitivité
\[
\underline{H}^*_Y(\Lambda_X) \Longleftarrow E_2^{pq} = \underline{H}^p_Y\bigl(\underline{H}^q_Z(\Lambda_X)\bigr) ,
\]
\struck{et l'hyp.\ de réc.} alors $\underline{H}^q_Z(\Lambda_X)$ est nul si $q \neq 2$, et loc.\ isom.\ à $\Lambda_Z$ si \uncertain{$r = 1$}, donc l'hyp.\ de réc.\ nous prouve que $E_2^{pq} = 0$ sauf pour \struck{$p = d$,} $p = 2(d-1)$, $q = 2$, et que pour ces valeurs de $p, q$, $E_2^{pq}$ est loc.\ isom.\ à $\Lambda_Z$, d'où la conclusion.\note{on attend « si $q = 2$ » là où la page semble porter « $r = 1$ ».}

Pour la définition d'un isom.\ canonique (4.2.1) ou (4.2.1 bis), le cas $d = 1$ est traité dans (SGA 4 XIX 1.3), \struck{le cas $d$} pour le cas $d$ quelconque, il suffit d'utiliser la construction de (SGA 4 XIX 3.4) ou de (SGA 4 XVIII \quad ).\note{« (4.2.1) ou (4.2.1 bis) » ainsi sur la page ; il s'agit de l'isomorphisme (4.2.2). La référence à SGA 4 XVIII est laissée en blanc.}

\page{104}
\note{Pagination de l'auteur : « 20 ».}

\textbf{Corollaire 4.3.} Soit $X$ un \struck{sous-}schéma \struck{\ill{}} régulier \struck{\ill{}} \add{satisfaisant au th.\ de pureté cf.\ (4.1)}, $Y$ un sous-schéma fermé régulier purement de codim.\ $d$, $n$ un entier premier aux caractéristiques résiduelles de $X$, $\Lambda = \mathbf{Z}/n\mathbf{Z}$, $F$ un faisceau loc.\ constant sur $X$ annulé par $n$. Pour tout faisceau de $\Lambda$-modules $G$ sur $X$ et tout $d \in \mathbf{Z}$, on pose $G(d) = G \otimes \mu_n^{\otimes(d)}$. \struck{Alors} On a une suite exacte canonique
\[
\begin{aligned}
(4.3.1) \qquad \cdots \to H^{i-2d}\bigl(Y, F(-d)\bigr) &\xrightarrow{\alpha} H^i(X, F) \xrightarrow{\beta} H^i(X - Y, F) \\
&\xrightarrow{\gamma} H^{i-2d+1}\bigl(Y, F(-d)\bigr) \to \cdots ,
\end{aligned}
\]
où $\beta$ est l'homomorphisme de restriction, et où l'homomorphisme composé
\[
H^{i-2}\bigl(X, F(-d)\bigr) \xrightarrow{\text{restr.}} H^i\bigl(Y, F(-d)\bigr) \xrightarrow{\alpha} H^i(X, F)
\]
est le cup-produit avec une classe de cohomologie\note{les exposants et les twists de cette ligne sont surchargés ; la lecture « $i-2$ », « $i$ » suit la page, non la cohérence avec (4.3.1).}
\[
(4.3.2) \qquad \gamma(Y) \in H^{2d}\bigl(X, \Lambda_X(d)\bigr) .
\]
\struck{Quand $Y$ \ill{} est un diviseur sur $X$,} Lorsque $d = 1$, cette classe est aussi la \add{première} classe de Chern de $\underline{L} = \underline{O}_X(Y)$, i.e.\ l'image par \add{l'homomorphisme cobord} \ill{} déduit de la suite exacte de Kummer
\[
0 \to \mu_{n\,X} \to \mathbf{G}_{m\,X} \xrightarrow{\lambda \mapsto \lambda^n} \mathbf{G}_{m\,X} \to 0
\]
\add{de la classe de $\underline{L}$ dans $\mathrm{Pic}(X) = H^1(X, \mathbf{G}_m)$}.

\emph{Démonstration.} \struck{Le} La suite exacte (4.3.1) n'est autre que la suite exacte de cohomologie relative
\[
\cdots \to H^i_Y(X, F) \to H^i(X, F) \to H^i(X - Y, F) \to \cdots ,
\]
compte tenu des isomorphismes
\[
H^i_Y(X, F) \simeq H^{i-2d}\bigl(Y, F(-d)\bigr)
\]
déduits de la suite spectrale de localisation
\[
H^*_Y(X, F) \Longleftarrow E_2^{pq} = H^p\bigl(Y, \underline{H}^q_Y(F)\bigr) ,
\]

\page{105}
\note{Pagination de l'auteur : « 21 ».}

et du calcul 4.2 des $\underline{H}^q_Y(F)$, impliquant que $E_2^{pq} = 0$ si $q \neq 2d$, d'où $H^i_Y(X, F) \simeq E_2^{i-2d,\,2d}$.

Faisant \struck{\ill{}} $F = \mu_n^{\otimes d}$ et $i = 2d$ dans (4.3.1), on trouve un homomorphisme
\[
H^{2d}_Y\bigl(X, \Lambda(d)\bigr) \simeq H^0(Y, \Lambda) \longrightarrow H^{2d}\bigl(X, \Lambda(d)\bigr)
\]
et l'image de la section \uncertain{unité} \add{constante $1$ de $\Lambda_Y$ sur $Y$} \struck{\ill{}} donne la \uncertain{première} classe (4.3.2). Les assertions \struck{\ill{}} de compatibilité \add{dans 4.3} \struck{faites} se prouvent par des arguments standard, pour lesquels nous renvoyons à SGA 5 IV.

\textbf{Remarque 4.3.3.} On \ill{} \ill{} \uncertain{seulement} utilisé l'hyp.\ de \ill{} \struck{\ill{}} avoir \ill{} les \ill{} de $\underline{H}^i_Y(\Lambda_X)$ ….\note{la remarque s'arrête là ; le reste du feuillet est blanc.}

\page{106}
\note{Pagination de l'auteur : « 22 ».}

\struck{4.3.\ \ill{} Soit $X$ un schéma régulier, et soit \ill{}} \struck{Rappelons \ill{} \ill{} \ill{} jusqu'en th.\ de} \struck{Abhyankar}\note{début biffé d'un numéro 4.3 bis, encadré et barré.}

\textbf{4.4.} Soit $X$ un schéma régulier, \struck{et soit} $(D_i)_{1 \leq i \leq r}$ une famille de \struck{diviseurs premiers} \struck{\ill{}} sous-schémas réguliers irréductibles dans $X$, tels que le diviseur $D = \sum D_i$ soit un diviseur à croisements normaux (SGA 5 I …). Soit $U = X - \operatorname{supp} D = X - \bigcup D_i$, $f : U \to X$ l'inclusion, \add{$n$ un entier premier aux} caractéristiques résiduelles de $X$ \struck{\ill{}} sur les points de \struck{supp} $\operatorname{supp} D$, $\Lambda = \mathbf{Z}/n\mathbf{Z}$. On se propose de calculer les faisceaux $R^i f_*(\Lambda_U)$ sur $X$. \struck{$R^i f_*(\Lambda_U)$}

\textbf{Proposition 4.5.} \struck{Avec les} \add{Sous les} conditions de 4.4, \struck{et \ill{} \ill{}} \add{supposons que $X$ satisfait au th.\ de pureté cohomologique (4.1)}\struck{$X$ satisfait au th.\ de pureté} :\note{l'hypothèse de pureté est écrite dans une bulle en marge gauche, reliée au texte.}

(i) \struck{$R^1 f_*(\Lambda_U)$ est canoniquement isom.} \add{On a un isom.\ canonique}
\[
(4.5.1) \qquad R^1 f_*(\Lambda_U) \simeq \coprod_{1 \leq i \leq r} \bigl(\mu_n^{\otimes -1}\bigr)_{D_i} \simeq \coprod_{1 \leq i \leq r} \Lambda_{D_i}(-1) .
\]

(ii) Pour tout entier $p$, l'homomorphisme \struck{canonique} \add{déduit du cup-produit}
\[
(4.5.2) \qquad \textstyle\bigwedge^p R^1 f_*(\Lambda_U) \longrightarrow R^p f_*(\Lambda_U)
\]
\add{sous l'hyp.\ de \uncertain{(ii)}}, est un isomorphisme.

(iii) Par suite, on a un isomorphisme canonique
\[
\begin{aligned}
(4.5.3) \qquad R^p f_*(\Lambda_U) &\simeq \coprod_{1 \leq i_1 < i_2 < \cdots < i_p \leq r} \bigl(\Lambda_X(-p)\bigr)_{D_{i_1 \ldots i_p}} \\
&\simeq \Lambda_X(-p) \otimes \coprod_{1 \leq i_1 < \cdots < i_p \leq r} \Lambda_{D_{i_1 \ldots i_p}} ,
\end{aligned}
\]
où $D_{i_1 \ldots i_p} = D_{i_1} \cap D_{i_2} \cap \cdots \cap D_{i_p}$.

\marginal{écrire $\Lambda_X(-p)$}
\note{dans (4.5.3), une première écriture biffée $(\mu_n^{\otimes -p})_X$, remplacée par $\Lambda_X(-p)$, comme l'indique la marge.}

\page{107}
\note{Pagination de l'auteur : « 23 ».}

\emph{Démonstration.} Soit $U_i = X - D_i$, $f_i : U_i \to X$ l'inclusion, \struck{de sorte que $f$ \ill{}} $f : U \to X$ se factorise en
\[
U \xrightarrow{\;g_i\;} U_i \xrightarrow{\;f_i\;} X , \qquad f = f_i g_i .
\]
On en conclut un homomorphisme canonique
\[
(*{*}) \qquad R^1 f_{i*}(\Lambda_{U_i}) \longrightarrow R^1 f_*(\Lambda_U) .
\]
D'autre part, \struck{(4.2.1) il \ill{}, \ill{} canonique} on a un isom.\ (4.2.2)
\[
R^1 f_{i*}(\Lambda_{U_i}) \simeq \underline{H}^2_{D_i}(\Lambda_X) \simeq \Lambda_{D_i}(-1) ,
\]
\note{après le second $\simeq$, une première écriture biffée $(\mu_n^{\otimes -1})_{D_i}$.}
d'où un homomorphisme
\[
\varphi_i : \Lambda_{D_i}(-1) \longrightarrow R^1 f_*(\Lambda_U) ,
\]
et par suite un homomorphisme
\[
\varphi : \coprod_{1 \leq i \leq r} \Lambda_{D_i}(-1) \longrightarrow R^1 f_*(\Lambda_U) .
\]
Je dis que c'est un homomorphisme \uncertain{et un isomorphisme}.

\note{la suite du feuillet 107 et le haut du feuillet 108, jusqu'à « … est un isomorphisme », forment un bloc encadré et barré par deux grandes diagonales ; il est lisible et transcrit ci-dessous.}

\struck{Pour cela, notons \ill{} que $U = \bigcup_{1 \leq i \leq r} U_i$, \uncertain{d'où} la suite spectrale de Leray}\note{« $\bigcup$ » ainsi sur la page, là où l'on attend une intersection.}
\[
\text{\struck{$R^* f_*(\Lambda_U) \Longleftarrow E_2^{pq} = H^p\bigl((i_0 \ldots i_p) \mapsto R^q f_{i_0 \ldots i_p *}(\Lambda_{U_{i_0 \ldots i_p}})\bigr)$}}
\]
\struck{où $U_{i_0 \ldots i_p} = U_{i_0} \cap \cdots \cap U_{i_p}$, et où $f_{i_0 \ldots i_p} : U_{i_0 \ldots i_p} \to X$ est l'inclusion canonique. On a évidemment}
\[
\text{\struck{$R^0 f_{i_0 \ldots i_p *}(\Lambda_{U_{i_0 \ldots i_p}}) = \Lambda_X$ ,}}
\]
\struck{\ill{} $E_2^{p0} = 0$ pour tout $p \neq 0$, d'où résulte que la suite spectrale donne un isomorphisme}
\[
\text{\struck{$(*{*}*) \quad R^1 f_*(\Lambda_U) \xrightarrow{\;\sim\;} E_2^{01} \subset \coprod_{1 \leq i \leq r} R^1 f_{i*}(\Lambda_{U_i})$}}
\]
\struck{Or on constate tout de suite que le}

\drawing{en marge gauche, petit schéma d'un quadrant $E_2^{pq}$ : deux axes, une croix à l'origine et une rangée de croix sur l'axe horizontal.}

\page{108}
\note{Pagination de l'auteur : « 24 ».}

\struck{composé de \uncertain{$\varphi$} avec l'hom.\ précédent est l'inclusion canonique des $R^1 f_{i*}(\Lambda_{U_i})$ dans la somme des $R^1 f_{i*}(\Lambda_{U_i})$, ce qui implique \ill{} (**). D'où la surjectivité de (**), et que $E_2^{01}$ est égal au dernier membre de (***), et que (**) est un isomorphisme. Cela prouve (i).} \struck{\ill{}} \struck{Supposons maintenant que $X$ satisfait au th.\ de pureté (ii). On} \struck{peut supposer} \add{est ramené à vérifier sur les fibres géométriques}, \add{donc on est ramené au cas} \struck{$X$ strictement local, et} \struck{de prouver alors que}
\[
\text{\struck{$\textstyle\bigwedge^p H^1(U, \Lambda) \longrightarrow H^p(U, \Lambda)$}}
\]
\struck{est un isomorphisme.}

\note{fin du bloc barré. Ce qui suit est la rédaction retenue.}

\uncertain{En} d'autres termes, \uncertain{il faut montrer} que l'hom.\ canonique
\[
(4.5.4) \qquad \coprod_{1 \leq i \leq r} (\Lambda)_{D_i} \longrightarrow R^1 f_*\bigl(\mu_{n\,X}\bigr)
\]
est un isomorphisme.\note{l'exposant de $\mu_{n\,X}$ est surchargé.} On \struck{peut} est, \add{ainsi que pour (ii)}, ramené pour ceci \struck{au cas où} à vérifier l'assertion fibre par fibre, ce qui nous ramène au cas où $X$ est strict.\ local, et à prouver que l'homom.\ canonique \uncertain{(ex)}
\[
\text{\struck{$\coprod$}} \quad \Lambda^r \longrightarrow H^1\bigl(U, \mu_n\bigr)
\]
\note{la ligne porte un $\coprod_{1 \leq i \leq r}$ biffé devant $\Lambda^r$ ; l'exposant de $H$ se lit « $1$ » ou « $r$ ».}
est un isom., ainsi que l'hom.
\[
\textstyle\bigwedge^i H^1(U, \Lambda) \longrightarrow H^i(U, \Lambda) ,
\]
dans le cas où $X$ satisfait au th.\ de pureté cohomologique.

\page{109}
\note{Pagination de l'auteur : « 25 ».}

La première assertion résulte \add{aussitôt} de la détermination du « groupe fondamental modéré » de $U$, rappelée plus bas (\quad ). Pour la deuxième, on procède par récurrence sur $r$, le cas $r = 0$ étant trivial, et le cas $r = 1$ étant contenu dans 4.2 (où on fait $d = 1$). Si $r \geq 2$, \struck{soit} considérons
\[
U' = X - \bigcup_{1 \leq i \leq r-1} D_i , \qquad D'_r = U' \cap D_r ,
\]
\note{devant $U' \cap D_r$, une première écriture biffée « $X \cap D_r$ ».}
de sorte que
\[
U = U' - D'_r , \qquad D'_r = D_r - \operatorname{supp} \Delta_r , \qquad \Delta_r = \sum_{1 \leq i \leq r-1} D_i \cap D_r .
\]
\note{après $U = U' - D'_r$, une ligne entière de formules lourdement biffée, illisible, se terminant par « $D_i \cap D_r$ ».}
\struck{\ill{}} Écrivons la suite exacte de cohomologie relative (4.3.1) pour $D'_r$ dans $U'$ :

\note{le reste du feuillet est barré de grandes diagonales ; c'est un premier état de la suite (4.5.5), repris au feuillet suivant. Transcrit tel quel.}

\struck{(4.5.5)} \quad $\cdots \to H^{i-2}(D'_r, \mu_n) \xrightarrow{\alpha_i} H^i(U', \Lambda) \to H^i(U, \Lambda) \to H^{i-1}(D'_r, \mu_n) \to \cdots$\note{le premier terme est surchargé ; une correction interlinéaire porte « $H^{i-2}(D'_r, \Lambda) \otimes \mu_n(k)^{-1}$ ».}

\struck{l'hyp.\ de récurrence \ill{} \ill{} \ill{} \ill{}, appliquée à $D'_r$} \add{suite spectrale}

\struck{D'autre part} la suite spectrale
\[
H^*_{D'_r}(U', \Lambda) \Longleftarrow E_2^{pq} = H^p\bigl(D'_r, \underline{H}^q_{D'_r}(\Lambda_{U'})\bigr)
\]
et la détermination 4.2 des $\underline{H}^q_{D'_r}(\Lambda_{U'})$ donne
\[
H^i_{D'_r}(U', \Lambda_{U'}) \simeq H^{i-2}\bigl(D'_r, \mu_n^{\otimes -1}{}_{D'_r}\bigr) \simeq H^{i-2}(D'_r, \Lambda) \otimes \mu_n^{\otimes -1}(k) ,
\]
\note{le dernier membre est biffé d'un trait léger.}
où $k$ est le corps résiduel de $X$. \struck{D'autre part l'hyp.\ de récurrence, appliquée à $D'_r$ et aux $r-1$ diviseurs \ill{}} D'autre part, l'hyp.\ de récurrence appliquée à $D'_r$ et aux $D_i \cap D_r$ ($1 \leq i \leq r-1$) sur $D'_r$ implique
\[
H^i(D'_r, \Lambda) \simeq \textstyle\bigwedge^i H^1(D'_r, \Lambda) \simeq \bigl(\bigwedge^i(\Lambda^{r-1})\bigr) \otimes \mu_n^{\otimes -i}(k)
\]
et de même l'hypothèse de récurrence appliquée à $U'$ et aux $D_i$ ($1 \leq i \leq r-1$) donne
\[
H^i(U', \Lambda) \simeq \textstyle\bigwedge^i H^1(U', \Lambda) \simeq \bigl(\bigwedge^i(\Lambda^{r-1})\bigr) \otimes \mu_n^{\otimes -i}(k) .
\]
\struck{Il faut exprimer \ill{} l'hom.\ de (*)}
\[
\text{\struck{$H^i_{D'_r}(U', \Lambda) \simeq H^{i-2}(D'_r, \Lambda) \otimes \mu_n^{\otimes -1}(k) \simeq \bigwedge^{i-2}(\Lambda^{r-1}) \otimes \mu_n^{\otimes((i-2)-1)}(k)$}}
\]
\[
H^i(U', \Lambda) \simeq \textstyle\bigwedge^i(\Lambda^{r-1}) \otimes \mu^{\otimes -i}(k)
\]
\note{les deux dernières lignes sont reliées par une boucle ; la dernière semble reprendre la précédente.}

\page{110}
\note{Pagination de l'auteur : « 26 ».}

\[
\begin{aligned}
(4.5.5) \qquad \cdots \to H^{i-2}\bigl(D'_r, \Lambda(-1)\bigr) &\xrightarrow{\alpha_i} H^i(U', \Lambda) \to H^i(U, \Lambda) \\
&\to H^{i-1}\bigl(D'_r, \Lambda(-1)\bigr) \to \cdots
\end{aligned}
\]
De plus, \struck{$H^*(U)$}, les termes $H^*(U', \Lambda)$ et $H^*(D'_r, \Lambda(-1))$ de cette suite exacte se calculent grâce à l'hypothèse de récurrence, et ils sont respectivement isomorphes à
\[
\textstyle\bigwedge^* H^1(U', \Lambda) \quad \text{et} \quad \bigl(\bigwedge^* H^1(D'_r, \Lambda)\bigr)(-1) .
\]
Je dis que\note{la phrase s'interrompt ici ; le reste du feuillet est blanc.}

\page{111}
\note{Pagination de l'auteur : « 27 ». Le feuillet continue la phrase « Je dis que » laissée en suspens au feuillet 110.}

\struck{Notons \ill{}} \struck{$H^{i-2}(D'_r, \Lambda) \otimes \mu_n(X)^{\otimes -1} \longrightarrow H^i$}\note{premier essai biffé, surchargé de traits.}

Je dis que l'hom.\ $\alpha_i$ de (4.5.5) est \emph{nul}. En effet, les calculs explicites précédents (liés à l'hyp.\ de récurrence) prouvent que
\[
H^*(U', F) \longrightarrow H^*(D'_r, F)
\]
est \struck{un isom.} \uncertain{surjectif} si $F$ est loc.\ constant et annulé par $n$, p.\ ex.\ si $F = \mu_n^{\otimes -1}$. Il suffit donc de prouver que le composé
\[
H^{i-2}\bigl(U', \Lambda(-1)\bigr) \longrightarrow H^{i-2}\bigl(D'_r, \Lambda(-1)\bigr) \xrightarrow{\alpha_i} H^i(U', \Lambda)
\]
est nul. Or on sait, \uncertain{d'après} 4.3, \struck{\ill{}} que le composé est le cup-produit par la classe \struck{du diviseur} $\gamma_{U'}(D'_r)$ \add{$= c^1(\underline{O}_{U'}(D'_r))$}, \struck{\ill{}} \uncertain{induite} par la classe \struck{\ill{}} $\gamma_X(D_r)$, qui est nulle puisque $X$ est str.\ local (donc que $D_r$ est principal). Donc la suite exacte (4.5.5) se décompose en suites exactes courtes
\[
(4.5.6) \qquad 0 \to H^i(U', \Lambda) \to H^i(U, \Lambda) \to H^{i-1}(D'_r, \Lambda)(-1) \to 0 .
\]
\note{le dernier terme porte un twist surchargé ; sous la suite, une formule biffée d'un trait ondulé, $\bigl(\bigwedge^{i-1}(\Lambda^{r-1})\bigr) \otimes \mu_n(k)^{\otimes -i}$, où l'on devine une première évaluation de ce terme.}

\struck{D'autre part, la \ill{} \ill{} $\bigwedge^* \bigl(\coprod_{1 \leq i \leq r} \Lambda_{D_i}\bigr)$ \ill{} isom.\ \ill{} $\Lambda^{r-1} \times \Lambda$, \ill{} \ill{} $\Lambda \to \Lambda_{D_r}$ \ill{}}\note{bas du feuillet encadré et biffé, en grande partie illisible.}

\page{112}
\note{Pagination de l'auteur : « 28 ».}

\struck{une suite exacte \ill{}}\note{en tête, une suite courte biffée $0 \to \bigwedge^i \; \bigwedge^{r-1} \otimes$, inachevée.}
\[
H^1(U, \Lambda) \simeq \coprod_{1 \leq i \leq r} H^1(U_i, \Lambda) \simeq \coprod_{1 \leq i \leq r-1} H^1(U_i, \Lambda) \times H^1(U_r, \Lambda)
\]
\[
\simeq H^1(U', \Lambda) \times \mu_n(X)^{-1}
\]
donc les \struck{suites exactes} \add{\uncertain{décompositions}}
\[
\textstyle\bigwedge^i H^1(U, \Lambda) \simeq \bigwedge^i H^1(U', \Lambda) \times \Bigl(\bigwedge^{i-1} H^1(U', \Lambda) \otimes \mu_n(X)^{-1}\Bigr)
\]
\struck{(4.5.7) \quad $0 \to \bigwedge^i H^1(U', \Lambda)$} qui définit une suite exacte
\[
(4.5.7) \qquad 0 \to \textstyle\bigwedge^i H^1(U', \Lambda) \longrightarrow \bigwedge^i H^1(U, \Lambda) \longrightarrow \Bigl(\bigwedge^{i-1} H^1(U', \Lambda)\Bigr)(-1) \to 0 .
\]
On \struck{\ill{}} envoie par cup-produit les termes de (4.5.7) dans les termes \struck{\ill{}} correspondants de (4.5.6). Nous admettrons qu'on obtient ainsi un diagramme de suites exactes (i.e.\ trivial pour le premier carré, et sensible dans le second, une vérification fastidieuse de \struck{\ill{}} compatibilité des \uncertain{isomorphismes} \struck{\ill{}}\note{la parenthèse ouverte avant « i.e. » se ferme au feuillet suivant, après la vérification ; quelques mots illisibles.}
\[
\text{\struck{$U = U' - D'_r$}}
\]
\add{\uncertain{l'hypothèse de récurrence}} \struck{\ill{}} explicite : \uncertain{$U'$ resp.\ $D'_r$}). \struck{\ill{}} \uncertain{Il s'ensuit}, \uncertain{puisque} les flèches verticales extrêmes \add{sont bijectives}, \struck{\ill{}} donc aussi la flèche verticale médiane, ce qui achève la démonstration.\note{l'ordre des lignes de ce bas de feuillet est incertain ; ajouts et ratures s'y chevauchent.}

\page{113}
\note{Pagination de l'auteur : « 29 ».}

\textbf{4.6.} Supposons, \struck{que} \struck{dans} sous les conditions de 4.4, que l'on ait pour tout $i$ ($1 \leq i \leq r$)
\[
D_i = V(a_i) , \quad \text{où} \quad a_i \in \Gamma(\underline{O}_X) .
\]
Soit $m$ un entier premier aux car.\ résiduelles de $X$ \struck{$\ill{}$ \ill{} $D_i$} et considérons le schéma
\[
X' = \operatorname{Spec} \underline{O}_X[T_1, \ldots, T_r]/(T_1^m - a_1, \ldots, T_r^m - a_r) .
\]
Je ne dis bien connu que $X'$ est un schéma régulier, et que \struck{\ill{}} si pour tout $i$ ($1 \leq i \leq r$), \struck{\ill{}} $b_i \in \Gamma(X', \underline{O}_{X'})$ est la section déduite de $T_i$, alors \add{les} $\Delta_i = V(b_i)$ sont des diviseurs réguliers de $X'$ tels que $\Delta = \sum \Delta_i$ soit un diviseur à croisements normaux (cf.\ SGA 1 XII …?). Donc on peut appliquer 4.5 à $X'$ et $U' = X' - \operatorname{supp} \Delta$. D'ailleurs, il est immédiat que $\operatorname{supp} \Delta_i$ \uncertain{est} \uncertain{l'image} inverse de $\operatorname{supp} D_i$, plus précisément, si $p : X' \to X$ est la projection,
\[
p^*(D_i) = m \Delta_i .
\]
On en conclut en particulier
\[
U' = p^{-1}(U) .
\]
Considérons l'homom.\ \struck{\ill{}} de changement de base
\[
(4.6.1) \qquad p^*\bigl(R^* f_*(\Lambda_U)\bigr) \longrightarrow R^* f'_*(\Lambda_{U'}) ,
\]
où $f' : U' \to X'$ est l'inclusion. \struck{\ill{}} On constate \add{\uncertain{alors}} aussitôt sur la définition des isomorphismes (4.5.1) que l'on a commutativité dans le diagramme

\page{114}
\note{Pagination de l'auteur : « 30 ».}

(4.6.2)

\begin{tikzcd}[column sep=large]
p^*\bigl(R^1 f_*(\Lambda_U)\bigr) \arrow[r, "{(4.6.1)}"] \arrow[d, "\simeq"'] & R^1 f'_*(\Lambda_{U'}) \arrow[dd, "\simeq"] \\
p^*\Bigl(\coprod_i \Lambda_{D_i}(-1)\Bigr) \arrow[d, "\simeq"'] & \\
\coprod_i \Lambda_{\Delta_i}(-1) \arrow[r, "m.\mathrm{id}"] & \coprod_i \Lambda_{\Delta_i}(-1)
\end{tikzcd}

où les flèches verticales \struck{\ill{}} correspondent aux isom.\ (4.5.1). Lorsque $X$ et $X'$ satisfont au th.\ de pureté, on en conclut de même la commutativité dans

(4.6.3)

\begin{tikzcd}[column sep=large]
p^*\bigl(R^p f_*(\Lambda_U)\bigr) \arrow[r, "{(4.6.1)}"] \arrow[d, "\simeq"'] & R^p f'_*(\Lambda_{U'}) \arrow[dd, "\simeq"] \\
p^* \bigwedge^p\Bigl(\coprod_i \Lambda_{D_i}(-1)\Bigr) \arrow[d, "\simeq"'] & \\
\bigwedge^p \coprod_i \Lambda_{\Delta_i}(-1) \arrow[r, "m^p.\mathrm{id}"] & \bigwedge^p \coprod_i \Lambda_{\Delta_i}(-1)
\end{tikzcd}

On en conclut en particulier :\note{la phrase se poursuit au feuillet suivant ; le reste du feuillet 114 est blanc.}

\page{115}
\note{Pagination de l'auteur : « 31 ».}

\textbf{Corollaire 4.7.} Avec les conditions de 4.4, supposons $X$ \add{$= \operatorname{Spec} A$} strict.\ local et satisfaisant aux \struck{conditions du} th.\ de pureté cohomologique. Soit $m$ un \struck{entier, \uncertain{premier}} multiple de $n$, premier à la car.\ résiduelle de $X$, et soit
\[
X'_m = \operatorname{Spec} A[T_1, \ldots, T_r]/(T_1^m - a_1, \ldots, T_r^m - a_r) ,
\]
où les $a_i$ sont des équations \struck{des} des $D_i$ ($1 \leq i \leq r$). Soit $U' = X' \times_X U$. Alors l'homomorphisme \uncertain{canonique}
\[
H^*(U, \Lambda) \longrightarrow H^*(U', \Lambda)
\]
est nul en dimension $> 0$.

En effet, c'est trivial en dim.\ $1$, en \add{\uncertain{utilisant}} (4.5.1) \struck{\ill{}} \struck{\ill{}} \add{\uncertain{équivaut}} précisément \add{à dire} que $U'$ est un revêtement \uncertain{principal} connexe, de groupe $\mu_n$.\note{phrase raturée et surchargée ; l'ordre des mots est incertain.}

\struck{4.8. Soit maintenant $X$ str.\ local \uncertain{régulier} \ill{} \ill{} th.\ de pureté cohomologique}\note{premier début de 4.8, encadré et barré.}

\textbf{4.8.} Soit $X$ \add{$= \operatorname{Spec} A$} un schéma strict.\ local régulier, $(a_i)_{1 \leq i \leq r}$ une partie d'un système régulier de paramètres de $A$, de sorte que \struck{$D = \operatorname{div}(a_i)$}
\[
D = \operatorname{div}\Bigl(\prod a_i\Bigr) = \sum D_i \qquad (\text{où } D_i = \operatorname{div} a_i)
\]
est un diviseur à croisements normaux, à composantes premières $D_i$ régulières. Posons
\[
U = X - \operatorname{supp} D = X - \bigcup D_i ,
\]
et pour tout entier $m$ premier à la car.\ résiduelle de $X$, introduisons
\[
X(m) = \operatorname{Spec} A[T_1, \ldots, T_r]/(T_1^m - a_1, \ldots, T_r^m - a_r) ,
\]

\page{116}
\note{Pagination de l'auteur : « 32 ».}

et l'image inverse $U(m)$ de $U$ dans $X(m)$, qui est un rev.\ étale principal de $U$ de groupe $\mu_m(X)^r = G(m)$, \struck{\ill{}} \add{\uncertain{formant}} un syst.\ projectif de \struck{schémas, \ill{}} revêtements étales principaux, \uncertain{compatibles} avec les homom.\ de transition \uncertain{évidents}
\[
\begin{array}{ccc}
\mu_{m'}^r(X) & \longrightarrow & \mu_m^r(X) \\
\| & & \| \\
G(m') & & G(m)
\end{array}
\qquad (m \mid m') .
\]
\note{au-dessus de « $(m \mid m')$ », le mot « pour » ajouté.}
Soit \struck{$U(\infty)$}
\[
U(\infty) = \varprojlim U(m)
\]
qui est un \struck{revêtement \uncertain{principal}} \uncertain{pro-revêtement} étale \uncertain{principal} de $U$ de groupe
\[
(4.8.1) \qquad G = \varprojlim_m G(m) = \prod_{\ell \neq p} \mathbf{Z}_\ell(1)^r ,
\]
($p$ car.\ la car.\ résiduelle de $X$, et où $\uncertain{(\ldots)}$ \ill{} \ill{} désigne)\note{la parenthèse est encerclée et partiellement illisible ; on attend la définition de $\mathbf{Z}_\ell(1)$ qui suit.}
\[
(4.8.2) \qquad \mathbf{Z}_\ell(1) = \varprojlim_\nu \mu_{\ell^\nu}(X) .
\]
Soit $n$ un entier $\geq 0$ premier à $p$, \uncertain{posons} $\Lambda = \mathbf{Z}/n\mathbf{Z}$. \add{et \uncertain{soit} \ill{}} On a
\[
H^*\bigl(U(\infty), \Lambda\bigr) \simeq \varinjlim_m H^*\bigl(U(m), \Lambda\bigr) ,
\]
\struck{et on peut \ill{} que}\note{encadré et barré :} \struck{$U(m) = X(m) - \bigcup_{1 \leq i \leq r} D'_i(m)$, avec $X(m)$ régulier et les $D'_i(m)$ $= \operatorname{div}(T_{i(m)})$ des diviseurs réguliers dont la somme est à croisements normaux, on trouve que} D'autre part, d'après les rappels de 4.6, le calcul de l'hom.\ de transition

\page{117}
\note{Pagination de l'auteur : « 33 ».}

\[
H^*\bigl(U(m), \Lambda\bigr) \longrightarrow H^*\bigl(U(m'), \Lambda\bigr) , \quad \text{pour } m \mid m'
\]
est justiciable de 4.6 \uncertain{qui} \struck{\ill{}} \add{\uncertain{en} degrés $\neq 0$} \uncertain{montre} \uncertain{qu'il} est nul si $nm \mid m'$, \uncertain{en} vertu du th.\ de pureté coh.\ \struck{D'ailleurs}, Cela prouve que la transition \uncertain{dans} (*) est nulle en dim.\ $> 0$, et en dim.\ $0$ il est \uncertain{évidemment} isomorphe à $\Lambda$. Donc on trouve

\textbf{Corollaire 4.9.} Avec les notations de 4.8, et supposons que $X$ et les $X(m)$ satisfont au th.\ de pureté coh., \struck{\ill{}} \ill{}. \uncertain{Alors}
\[
(4.9.1) \qquad
\begin{cases}
H^i\bigl(U(\infty), \Lambda\bigr) = 0 & \text{si } i \neq 0 \\
H^0\bigl(U(\infty), \Lambda\bigr) \simeq \Lambda .
\end{cases}
\]
\struck{\textbf{Remarques 4.10.} a) \ill{} \ill{} \uncertain{seulement} que les $X(m)$ \uncertain{satisfont} au th.\ de pureté \uncertain{coh.}\ en \uncertain{dim.}\ $\leq N$, \ill{} \uncertain{condition} \ill{} \uncertain{toujours} \uncertain{vérifiée} si $N = 1$ !), alors (4.9.1) \uncertain{reste} vrai \uncertain{pour} $i \leq N$.}

\struck{b) On \uncertain{exprime} \uncertain{encore} \ill{} \ill{} résultat de 4.9 en disant que \ill{} \ill{} \ill{} \uncertain{fois} \ill{} \ill{} $U$ \ill{} \ill{} \uncertain{les} \ill{} \ill{} \uncertain{premier} \ill{}}\note{les remarques 4.10 a) et b) sont encadrées et barrées de grandes diagonales ; le passage est en grande partie illisible.}

\textbf{Corollaire 4.10.} \add{Sous les conditions de 4.9,} Soit $\Gamma$ un groupe profini, considérons un homomorphisme de groupes profinis
\[
(4.10.1) \qquad \varphi : G \longrightarrow \Gamma ,
\]
et soit $H$ le noyau de cet homomorphisme.\note{la lettre avant « $\longrightarrow \Gamma$ » est surchargée ; lue $G$, le groupe de (4.8.1).}

\page{118}
\note{Pagination de l'auteur : « 34 ». Sur ce feuillet et le suivant, le groupe de (4.8.1) est écrit « $G$ » suivi d'un griffonnage (peut-être « $G(\infty)$ ») ; on transcrit $G$. Le signe $\overset{G}{\times}$ est le produit contracté.}

\struck{Alors} Considérons l'anneau \add{\uncertain{de cohomologie}} $H^*(H, \Lambda)$, sur lequel $G$ \uncertain{opère} trivialement, et considérons l'anneau induit
\[
H^*(H, \Lambda) \overset{G}{\times} \Gamma .
\]
\struck{Alors} Considérons d'autre part le revêtement principal $\widetilde{U}'$ de groupe $\Gamma$
\[
\widetilde{U}' = U(\infty) \overset{G}{\times} \Gamma
\]
associé à $U(\infty)$ de groupe $G$, et à l'hom.\ $G \to \Gamma$. \uncertain{Alors} \uncertain{on a} \uncertain{donc} un isomorphisme canonique \add{\uncertain{gradué} à opérateurs $\Gamma$} d'anneaux\note{une boucle relie ce passage à la formule (4.10.1) ; une phrase intermédiaire, « Faisant opérer $\Gamma$ sur l'anneau $H^*(\widetilde{U}, \Lambda)$ via ses opérations sur $\widetilde{U}$ », est encerclée et rattachée par une flèche.}
\[
(4.10.1) \qquad H^*(\widetilde{U}', \Lambda) \simeq H^*(H, \Lambda) \overset{G}{\times} \Gamma .
\]
\struck{En particulier, \ill{} $\Gamma'$ \ill{} \ill{} $\Gamma' \subset \Gamma$ \ill{} $\Gamma'$ \ill{} \ill{} $G$ \ill{} $\Gamma$, \ill{} \ill{} conséquence \ill{} triviale}\note{deux lignes barrées, très surchargées.}

\emph{Démonstration.} Soit $K = G/H = \operatorname{Im}(G \to \Gamma)$. Alors
\[
U' = (U(\infty)/H) \overset{K}{\times} \Gamma ,
\]
et par suite on trouve un isom.\ d'anneaux gradués à \uncertain{groupe} d'opér.\ $\Gamma$
\[
H^*(U', \Lambda) \simeq H^*(U(\infty)/H, \Lambda) \overset{K}{\times} \Gamma
\]
où $K$ opère sur $H^*(U(\infty)/H, \Lambda)$ via ses opérations sur $U(\infty)/H$. \uncertain{Il} \uncertain{nous} \uncertain{suffit} \uncertain{donc} de \uncertain{montrer} que en dimension $\neq 0$ \struck{\ill{}} \add{\uncertain{triviales}},\note{lecture incertaine de cette ligne ; on attend que $K$ opère trivialement.} \struck{\ill{} \ill{}} et qu'on a un isom.\ d'anneaux gradués
\[
H^*(U(\infty)/H, \Lambda) \simeq H^*(H, \Lambda) ,
\]
\uncertain{Cela} \uncertain{nous} \uncertain{ramène} au cas \uncertain{où} $\varphi$ \uncertain{est} \uncertain{supposé} \emph{surjectif}. Alors la \uncertain{suite} \uncertain{spectrale} \ill{} (4.10.1) \uncertain{est} la suite spectrale dite de Hochschild-Serre (SGA 4 VII …) de
\[
H^*(U(\infty)/H, \Lambda) \Longleftarrow E_2^{pq} = H^p\bigl(H, H^q(U(\infty), \Lambda)\bigr) ,
\]
\note{l'inégalité « en dimension $\neq 0$ » et la structure de la phrase précédente sont reconstruites d'après la page ; « SGA 4 VII » porte un trait au-dessus du chiffre romain.}

\page{119}
\note{Pagination de l'auteur : « 35 ».}

\uncertain{jointe} à 4.9, qui indique que $E_2^{pq} = 0$ si $q \neq 0$, donne l'isomorphisme voulu d'anneaux gradués :
\[
H^*(U(\infty)/H, \Lambda) \simeq H^*(H, \Lambda) .
\]
D'autre part, cet hom.\ est, \uncertain{comme} \uncertain{bien} connu, \uncertain{transforme}, \uncertain{comme} \ill{}, il est \uncertain{compatible} aux opérations de $G$ sur les \uncertain{premiers} \uncertain{membres} \uncertain{(}les \uncertain{quotients} \ill{} des opérations de $G$ sur $H$ par automorphismes intérieurs. Comme $G$ est commutatif, cette action est triviale, ce qui achève la démonstration de 4.10.\note{deux lignes du milieu de phrase lues avec peine ; le sens attendu est que $G$ opère sur $H^*(H, \Lambda)$ par automorphismes intérieurs, donc trivialement, $G$ étant commutatif.}

\textbf{Remarques 4.12.} \uncertain{Si} \uncertain{on} \uncertain{suppose} seulement que les $X(m)$ satisfont au th.\ de pureté coh.\ en dim.\ $\leq N$, \uncertain{alors} les arguments qui précèdent \uncertain{prouvent} encore 4.9 à 4.11 en dim.\ $\leq N$, \struck{et} \struck{\ill{} corollaire 4.13 qui suit}.\note{une flèche renvoie le numéro « 4.12 » d'un en-tête biffé plus bas jusqu'à cette remarque, qui était d'abord numérotée autrement.}

\textbf{Corollaire 4.11.} \add{Soit $\Gamma'$ l'image de $G$ dans $\Gamma$, et} \uncertain{supposons} $\Gamma$ commutatif. \struck{On \ill{} $\Gamma'$ invariant $\Gamma$ \ill{} \ill{} \ill{} $G(\infty)$ dans $\Gamma$}.\note{lecture du début incertaine ; « $G(\infty)$ » est ici écrit en toutes lettres, dans un passage biffé.} Alors, \struck{4.12.} plus gén.\ \struck{\ill{}} \struck{\ill{}}, l'action de $\Gamma$ sur $H^*(U', \Lambda)$ induit sur $\Gamma''$ l'action triviale \add{(plus précisément $\Gamma''$ est le s.-groupe des $g \in \Gamma$ qui opèrent trivialement sur $H^*(U', \Lambda)$)}.\note{le passage est encadré à gauche ; la parenthèse est une insertion interlinéaire. $\Gamma''$ n'est pas défini dans ce qui se lit.}

\uncertain{Cet} \uncertain{énoncé} \uncertain{simplifie} \struck{\ill{}} \uncertain{beaucoup} \ill{} \struck{\ill{} maintenant}

\struck{4.13. Nous \uncertain{prouverons}} \struck{plus bas} plus bas le th.\ de \uncertain{monodromie} \uncertain{sous} la forme 3.2. Avant de le faire, explicitons le calcul de $H^*(U', \Lambda)$ résultant de (4.10.1), ce qui nous sera utile dans l'exposé suivant.\note{« 3.2 » renvoie à un énoncé antérieur au lot, que ces pages ne permettent pas d'identifier.}

\page{120}
\note{Pagination de l'auteur : « 36 ».}

\textbf{4.13.} Gardons les hypothèses de 4.9. \struck{\ill{} \ill{}} \struck{\ill{} \ill{} \ill{} la forme $\ell^{\nu+1}$ \ill{}}
\[
(4.13.1) \qquad \Lambda = \mathbf{Z}/\ell^{\nu+1}\mathbf{Z} ,
\]
\note{la ligne (4.13.1) est encadrée et rattachée à la ligne biffée qui la précède.}
Le groupe commutatif $H$ se décompose en le produit de ses composantes $\ell$-primaires $H(\ell)$,
\[
H = \prod_{\ell \neq p} H(\ell) , \qquad H(\ell) \subset G(\infty)(\ell) \simeq \mathbf{Z}_\ell(1)^r .
\]
\note{devant $\mathbf{Z}_\ell(1)^r$, un signe biffé, peut-être un $\prod$.}
Donc $H(\ell)$ est un \emph{module libre} sur $\mathbf{Z}_\ell$ de rang $s \leq r$. \struck{$\Lambda = \mathbf{Z}/n\mathbf{Z}$ avec $n = \ell^{\nu+1}$, \ill{} $\Lambda$ \ill{} un \ill{} \ill{}}

On en conclut \struck{\uncertain{par}} aussitôt, si
\[
(4.13.2) \qquad \Lambda = \mathbf{Z}/n\mathbf{Z} , \quad \text{avec } n = \ell^{\nu+1} ,
\]
\struck{\uncertain{puissance} \ill{} de $\ell$}
\[
(4.13.1) \qquad
\begin{cases}
H^*(H, \Lambda) \simeq \bigwedge^* H^1(H, \Lambda) \\
H^1(H, \Lambda) \simeq \operatorname{Hom}(H, \Lambda) \simeq \Lambda^s , \quad \text{où } H(\ell) \simeq \mathbf{Z}_\ell^s .
\end{cases}
\]
\note{le numéro « (4.13.1) » est répété sur la page ; un mot biffé illisible suit $\operatorname{Hom}$ dans la première ligne.}

\textbf{4.14.} \struck{Soit maintenant} \struck{\ill{}} Supposons \struck{maintenant} \uncertain{maintenant} donnés des entiers $m_i \geq 0$ ($1 \leq i \leq r$), et \struck{considérons} soit
\[
(4.14.1) \qquad t = a_1^{m_1} \cdots a_r^{m_r}
\]
\note{devant $t$, un signe biffé.}
Soit, pour tout entier $m$ premier à la car.\ résiduelle de $X$,
\[
(4.14.2) \qquad X'(m) = \operatorname{Spec} A[T]/(T^m - t)
\]
et \struck{considérons}
\[
(4.14.3) \qquad \text{\struck{$X(m) =$}} \quad U'(m) = X'(m) \times_X U ,
\]
qui est un revêtement étale principal de $U$, de groupe $\mu_m(X)$. Posons \add{$Y'(m)$, et les} $U'(m)$, formant un système projectif, et \uncertain{passant} à la limite
\[
(4.14.4) \qquad U' = \varprojlim_m U'(m) ,
\]
qui est un schéma \uncertain{principal} \uncertain{sur} $U$ de groupe
\[
(4.14.5) \qquad \Gamma = \varprojlim \mu_m(X) = \prod_{\ell \neq p} \mathbf{Z}_\ell(1) .
\]
\note{la phrase « Posons … » est corrigée en interligne et se lit mal ; la page s'arrête sur (4.14.5), et l'argument se poursuit au-delà du lot. Le « $\ell \neq 1$ » apparent sous le dernier produit est lu $\ell \neq p$.}

\end{document}
